\documentclass[11pt,a4paper]{article}

\usepackage{main}
\usepackage{cours}

\renewcommand{\typedocument}{Cours}
\renewcommand{\classe}{3e}
\renewcommand{\titre}{Mouvement}

\begin{document}

\pagination

\creertitre{Mouvement et vitesse}

\partie{Définitions}

On appelle \important{mouvement} le déplacement d’un objet par rapport à un autre. Un objet est dit en mouvement lorsque sa position varie au cours du temps.

On appelle \important{mobile} l’objet en mouvement. Il est souvent représenté par un \textbf{point} dans les exercices.

On appelle \important{référentiel} l’objet par rapport auquel on décrit le mouvement. Il s’agit de l’observateur du mouvement.

On appelle \important{trajectoire} l’ensemble des positions prises par un mobile au cours du temps.

\partie{Les référentiels}

Le mouvement d’un mobile dépend du référentiel choisi.

\exemple{Un passager assis dans un train est immobile par rapport au train, mais il sera en mouvement par rapport au sol.}

Selon la situation, certains référentiels sont plus adaptés :

\souspartie{Le référentiel terrestre local}

Le \important{référentiel terrestre local} est lié à la surface de la Terre, le “sol”.

On l’utilise le plus souvent pour décrire les mouvements du quotidien.

\exemple{une voiture sur une route ; un ballon lancé dans la cour.}

\souspartie{Le référentiel géocentrique}

Le \important{référentiel géocentrique} a pour origine le centre de la Terre.

On l’utilise pour étudier le mouvement d’objets situés autour de la Terre.

\exemple{la Lune autour de la Terre ; des satellites artificiels.}

\souspartie{Le référentiel héliocentrique}

Le \important{référentiel héliocentrique} a pour origine le centre du Soleil.

On l’utilise pour étudier le mouvement des planètes dans le Système solaire.

\exemple{la Terre, Mars, etc...}

\partie{Les types de mouvements}

Si la trajectoire suit une droite, le mouvement est dit \important{rectiligne}. 

\exemple{Une pomme qui tombe d’un arbre}

Si la trajectoire suit un arc de cercle, le mouvement est dit \important{circulaire}.

\exemple{La Terre autour du Soleil}

Si la trajectoire suit une courbe, le mouvement est dit \important{curviligne}.

\exemple{La courbe d’un boulet de canon}

\newpage

\partie{La vitesse}

Pour n’importe quel objet en mouvement d’un point A à un point B en un temps T, on peut calculer la \important{vitesse moyenne}:

\important{v = d/T} où d est la distance entre A et B et T le temps du trajet.

La vitesse d’un mobile à un instant donné s’appelle la \important{vitesse instantanée}. Elle est égale à la vitesse moyenne uniquement pour un mouvement uniforme.

\exemple{Lors d'un trajet en voiture on ne roule pas toujours à la même vitesse, la vitesse instantanée (affichée au compteur) n'est donc pas tout le temps égale à la vitesse moyenne (distance totale parcourue / temps total du trajet).}

\souspartie{Variation de vitesse et mouvement}

Si la vitesse d’un mobile \textbf{augmente}, le mouvement est \important{accéléré}.

Si la vitesse d’un mobile \textbf{diminue}, le mouvement est \important{ralenti}.

Si la vitesse d’un mobile est \textbf{constante}, le mouvement est \important{uniforme}.

\souspartie{Représentation de la vitesse}

Sur un schéma, on représente la vitesse par une flèche dont :

\begin{itemize}
    \item la \textbf{direction} est celle de la trajectoire,
    \item le \textbf{sens} est le même que celui du mouvement,
    \item la \textbf{longueur} est proportionnelle à la valeur de la vitesse.
\end{itemize}

\end{document}